\section{Introduction}
\label{sec:introduction}

Retrieval-augmented generation over PDFs first turns each document into
text and cuts the text into chunks; a retriever then picks the chunks a
language model reads. How to cut is an open choice, and the number of
chunking methods keeps growing
\citep{smigielski2026chunking,shaukat2026systematic,zhao2025moc,lu2025hichunk}.
Because no method wins everywhere, recent work proposes choosing the
chunking method per document \citep{demourajunior2026adaptive}. Whether
that choice pays off is hard to judge from published comparisons: they
usually retrieve a fixed number of chunks, so a method with larger chunks
also gives the reader more text, and the choice is usually tuned and
scored on the same questions.

We ask five questions. \textbf{RQ1}: does choosing a setup (parser,
chunker and chunker setting) per document beat one fixed setup?
\textbf{RQ2}: how do chunk size and chunker affect retrieval at a fixed
token budget? \textbf{RQ3}: does the measure change which size looks
best? \textbf{RQ4}: do answers from a language model follow retrieval?
\textbf{RQ5}: do tuned chunk sizes and parser settings beat their
defaults?

To answer RQ1 fairly, we score every setup by the share of gold-answer
characters read within the same token budget, so chunk size and retrieval
depth trade off, and we estimate per-document choice with a split-half
oracle that chooses on half of a document's questions and scores on the
other half. The study was pre-registered: the scorer, the analysis and
the decision rule were frozen and pushed to a public repository before
the held-out documents were opened, and every later change is reported
as a dated deviation.

Our contributions are:
\begin{itemize}
\item a pre-registered comparison of per-document choice against one
fixed setup, with \RQoneSetups{} setups on \AccountingMain{} held-out
documents, two parsers and a frozen scorer;
\item an evaluation at an equal token budget, with a structure-tree
chunker of our own among the candidates;
\item evidence that the measure decides which chunk size looks best, and
that a library default costs more than per-document choice gains;
\item the plan, deviations, scorer and result files: every number in
this paper is generated from a checksummed result file or, where noted,
taken from a dated specification note.
\end{itemize}

The answer to RQ1 is \RQoneDecision{}: the per-document oracle does not
beat the best single setup (gap \RQonePrimaryGap{}
\RQonePrimaryGapCI{}), and the recommendation is one fixed setup with a
small, tuned chunk size.
